\documentclass[11pt]{article}
\usepackage[T1]{fontenc}
\usepackage{lmodern}
\usepackage{amsmath,amssymb,amsthm}
\usepackage[margin=1in]{geometry}
\usepackage{booktabs}
\usepackage{microtype}
\usepackage[colorlinks=true,allcolors=blue]{hyperref}
\usepackage{enumitem}
\newtheorem{theorem}{Theorem}
\newtheorem{lemma}[theorem]{Lemma}
\newtheorem{corollary}[theorem]{Corollary}
\newtheorem{proposition}[theorem]{Proposition}
\theoremstyle{remark}\newtheorem{remark}[theorem]{Remark}
\newcommand{\MT}{\mathop{\rm MinTr}}
\newcommand{\poly}{\mathop{\rm poly}}
\newcommand{\ind}[1]{\mathbf 1\{#1\}}
\newcommand{\sat}{\#\mathrm{SAT}}
\newcommand{\eps}{\varnothing}
\title{Linear-variable exact counting of minimal transversals\\
\large A completed reduction, size refinement, and algorithmic consequences}
\author{}
\date{29 October 2026}
\begin{document}
\maketitle
\begin{abstract}
For an explicitly represented hypergraph with $p$ vertices and $q$ indexed
edges, we give a polynomial-time reduction of its number of inclusion-minimal
transversals to the difference of two ordinary width-at-most-three CNF counts.
When $p+q\ge1$, each formula has exactly $3p+2q-1$ declared variables in the
unoptimized construction. The reduction has no auxiliary variable per incidence
or per pair of vertices in an edge. A uniquely extendible linear-size population
count circuit also gives the entire cardinality distribution using polynomially
many calls. Closed-neighborhood hypergraphs yield $5n-1$ variables for counting
minimal dominating sets of arbitrary $n$-vertex graphs. Combining the reduction
with the cited fixed-width SAT counting theorem gives deterministic
subexponential time and polynomial space. The reduction and its proof are
independent of that external theorem. The cancellation mechanism has prior
antecedents in negative-weight model counting; no priority claim is made.
\end{abstract}

\section{Scope and main statement}
This note studies minimal transversal counting and its
domination and set-cover applications, including size
refinement and witness recovery. The optional connectivity and feedback vertex
set directions are discussed separately in Section~\ref{sec:limits}; they are
not consequences of the private-edge theorem here. The accompanying note
\texttt{papers/connected\_domination.tex} establishes a separate randomized
decision and search result for connected domination using a monomial-weight
extension of the external counting algorithm. Exact connected-domination
counting and the feedback-vertex-set direction remain open in this project.

Let $H=(U,(e_i)_{i=1}^q)$, where $U=\{1,\ldots,p\}$ and $e_i\subseteq U$.
Repeated edges are permitted and retain distinct indices. A transversal $D$
meets every edge. It is minimal if no proper subset is a transversal. Write
\[
 s=p+q,\qquad I=\sum_{i=1}^q |e_i|,\qquad
 J=\sum_{i=1}^q\binom{|e_i|}{2}.
\]
Input size includes the declared vertices, edge indices, and their incidence
lists. All polynomial bounds below refer to this explicit representation.

\begin{theorem}[Two-call reduction]\label{thm:main}
For $s\ge1$, one can construct two CNF formulas $F_0,F_1$, each of width at most
three, such that
\[
 |\MT(H)|=\sat(F_0)-\sat(F_1).
\]
Before any simplification, each formula has
\[
 N=3p+2q-1,\qquad M=p+I+J+4(s-1)+1
\]
declared variables and clauses, respectively. The construction takes polynomial
time and space. Fixing any subset of the $p$ selection variables gives the same
identity for minimal transversals satisfying those fixed choices.
For $p=q=0$, the answer is $1$ and may instead be returned directly.
\end{theorem}

This is a two-query difference reduction, not a
parsimonious reduction, not projected model counting, and not a parity-only
algorithm. In particular, existence cannot in general be detected merely by
testing the satisfiability of either formula.

\section{The signed identity}
For $D\subseteq U$, define the set of its private-edge owners
\[
 W(D)=\{v\in D:\text{there is an }i\text{ with }e_i\cap D=\{v\}\}.
\]
The familiar private-edge characterization can be proved directly: if $D$ is
a transversal, removing $v$ fails precisely when some edge meets $D$ only in
$v$. If every one-element deletion fails, every proper-subset deletion fails
as well, because the transversal property is upward closed. Consequently,
\begin{equation}\label{eq:private}
 D\in\MT(H)\quad\Longleftrightarrow\quad
 D\text{ is a transversal and }W(D)=D.
\end{equation}

Introduce three sets of Boolean choices:
\[
 D,R\subseteq U,\qquad P\subseteq\{1,\ldots,q\}.
\]
Call $(D,R,P)$ admissible when
\begin{equation}\label{eq:admissible}
 R\subseteq D,\qquad
 i\in P\ \Longrightarrow\ |e_i\cap D|\le1\text{ and }e_i\cap D\subseteq R.
\end{equation}
Associate the sign $(-1)^{|R|+|P|}$.

\begin{lemma}[Pointwise cancellation]\label{lem:pointwise}
For every $D\subseteq U$,
\begin{equation}\label{eq:pointwise}
 \sum_{\substack{R,P\\(D,R,P)\text{ admissible}}}
 (-1)^{|R|+|P|}=\ind{D\in\MT(H)}.
\end{equation}
\end{lemma}
\begin{proof}
Fix $D$ and $R\subseteq D$. The eligible edge-index set is
\[
 A(D,R)=\{i:|e_i\cap D|\le1,\ e_i\cap D\subseteq R\}.
\]
Summing over its arbitrary subsets gives
\[
 \sum_{P\subseteq A(D,R)}(-1)^{|P|}
 =\begin{cases}1&A(D,R)=\eps,\\0&A(D,R)\ne\eps.\end{cases}
\]
If $D$ misses an edge $e_i$, that index belongs to $A(D,R)$ for every $R$;
the whole contribution is therefore zero. Suppose instead that $D$ is a
transversal. Every eligible edge has a unique intersection vertex in $R$.
Thus $A(D,R)=\eps$ exactly when $R\cap W(D)=\eps$. It follows that the left
side of \eqref{eq:pointwise} is
\[
 \sum_{R\subseteq D\setminus W(D)}(-1)^{|R|}
 =\begin{cases}1&D\setminus W(D)=\eps,\\0&\text{otherwise}.\end{cases}
\]
Equation~\eqref{eq:private} finishes the proof, including $D=\eps$.
\end{proof}

For example, take $U=\{1,2\}$ and the single edge $\{1,2\}$. The contributions
of $D=\eps,\{1\},\{2\},\{1,2\}$ are respectively $0,1,1,0$. The signed sum
is $2$, regardless of how many identical copies of that edge are indexed.

\subsection{Simplifying the earlier four-set construction}
The earlier four-set construction used $D,T,P,Q$ and sign
$(-1)^{|D|+|T|+|P|+|Q|}$. Its separate $Q$ filter for missed edges is unnecessary:
a missed edge already makes the corresponding $P$ bit free, hence cancels.
Removing $Q$ leaves a valid signed sum. Substituting $R=D\setminus T$ gives
\eqref{eq:admissible}, and
\[
 |D|+|T|\equiv |R|\pmod2.
\]
This also removes the $D$ bits from the parity circuit. The improvement is in
variable count; many clauses are still permitted.

\section{Explicit width-three formulas}
Use variables $x_v$ for $v\in D$, $r_v$ for $v\in R$, and $a_i$ for $i\in P$.
The base formula $B_H$ consists of the following clauses:
\begin{align}
 &\neg r_v\vee x_v &&(v\in U),\label{eq:c1}\\
 &\neg a_i\vee\neg x_v\vee r_v &&(v\in e_i),\label{eq:c2}\\
 &\neg a_i\vee\neg x_v\vee\neg x_w
 &&(\{v,w\}\subseteq e_i,\ v\ne w).\label{eq:c3}
\end{align}
Its $2p+q$ variables and $p+I+J$ clauses express precisely
\eqref{eq:admissible}. In \eqref{eq:c3} each unordered pair is used once.
The construction does not introduce a variable for an incidence or a pair.

Order the $s$ variables in $(r_1,\ldots,r_p,a_1,\ldots,a_q)$ as
$b_1,\ldots,b_s$. Compute their parity using $s-1$ successive XOR gates, with
the first bit as the initial wire. Encode an XOR $z=u\mathbin\oplus v$ by
\begin{equation}\label{eq:xor}
 (u\vee v\vee\neg z)\wedge
 (u\vee\neg v\vee z)\wedge
 (\neg u\vee v\vee z)\wedge
 (\neg u\vee\neg v\vee\neg z).
\end{equation}
Conjoin a unit clause fixing the final parity wire to $j$ to obtain $F_j$.
For $s=1$, that wire is $b_1$ itself. Every base assignment has exactly one
extension to the XOR variables, so the difference of the two model counts
equals the signed sum of Lemma~\ref{lem:pointwise}. The variable and clause
counts in Theorem~\ref{thm:main} follow. Additional unit clauses on $x$ restrict
the outer sum over $D$, establishing the theorem's conditioning assertion.

All variables must remain declared in a DIMACS export, even when unused in
the base constraints. In particular, an empty edge has a free $a_i$ before
parity is imposed; omitting it would destroy cancellation. Width less than
three is allowed by the imported counting theorem. If exactly three literal
positions are required by an interface, repeating a literal does not change
any model count.

\subsection{Boundary cases}
An empty edge forces the answer to zero and is also handled by the formula.
When $q=0$, the unique minimal transversal is $\eps$, including if $p>0$.
When $p=0<q$, every edge is empty and the answer is zero. Repeated edges do
not affect the proof; they merely create more eligible indices in a cancelling
sum. An isolated vertex cannot belong to a minimal transversal: it never lies
in $W(D)$. Isolated graph vertices are different, since their closed
neighborhoods are singleton edges and must be selected in a dominating set.

\section{Cardinality distributions with linear overhead}\label{sec:size}
Define
\[
 m_k(H)=|\{D\in\MT(H):|D|=k\}|,\qquad
 M_H(z)=\sum_{k=0}^p m_k(H)z^k.
\]
The pointwise lemma implies the polynomial identity
\[
 M_H(z)=\sum_{(D,R,P)\text{ admissible}}(-1)^{|R|+|P|}z^{|D|}.
\]
To extract a coefficient, conjoin a deterministic circuit testing
$\sum_v x_v=k$. A generic sequential binary accumulator would cost
$O(p\log p)$ gates; the following balanced construction costs $O(p)$.

For $p\ge1$, pad the input with constant zeros to $b=2^h$, the smallest power
of two at least $p$. At level $j=1,\ldots,h$, use $b/2^j$ additions of two
$j$-bit numbers. The low bit uses one half adder; each remaining bit uses a
full adder. A half adder needs two output variables (XOR and AND) and seven
clauses. A full adder uses two XOR outputs for its sum and a majority output
for its carry: three variables and fourteen clauses. A carry $c$ for input
bits $u,v,w$ is enforced by
\[
 \bigwedge_{\{\alpha,\beta\}\subseteq\{u,v,w\}}
 [ (\neg\alpha\vee\neg\beta\vee c)
   \wedge(\alpha\vee\beta\vee\neg c) ].
\]
All clauses have width at most three and uniquely determine the output bits.
Constants and aliases can be simplified without adding variables. An
unsimplified upper bound on new variables is
\begin{equation}\label{eq:adder}
 A_p\le\sum_{j=1}^h\frac b{2^j}(3j-1)
 =5b-3h-5<10p,
\end{equation}
and on gate clauses it is
\[
 \sum_{j=1}^h\frac b{2^j}(14j-7)=21b-14h-21<42p.
\]
Finally fix the $h+1$ output bits to the binary expansion of $k$. For $p=1$
there are no additions; for $p=0$ the constraint is constant. An out-of-range
$k$ gives count zero.

\begin{corollary}\label{cor:size}
For every $k$, $m_k(H)$ is the difference of two width-three counts on
$3p+2q-1+O(p)$ variables. All $p+1$ coefficients are obtainable with
$2(p+1)$ oracle calls (apart from trivial cases). Every gate extension is
unique, so there is no unintended witness multiplicity.
\end{corollary}

Threshold comparison with $k$ also has $O(\log(p+2))$ extra gates after the
same population count. Thus maximum feasible cardinality can alternatively
be found with $O(\log(p+2))$ count-difference tests of the monotone predicate
``there exists a minimal transversal of size at least $k$''.

\section{Consequences of the external SAT theorem}\label{sec:time}
We use the October 25 manuscript's fixed-width counting theorem as an
external algorithmic input~\cite{oai}. Its main theorem gives, for fixed
width, time
\[
 (N+M)^{O(1)}2^{O(N\log\log N/\sqrt{\log N})}
\]
and polynomial space. For $M=O(N)$ the exponential factor improves to
$2^{O(N/\sqrt{\log N})}$, with constants depending on the density bound.
The linked Lean scope document~\cite{scope} states the weaker factor
$2^{O(N(\log\log N)^2/\sqrt{\log N})}$ for canonical encodings, also with
polynomial space. We have inspected the stated theorem and scope; this work
does not rebuild that formalization or reprove its algorithm. The sharper
bound below uses the manuscript, not a claim that the sharper exponent was
Lean-verified.

For small arguments, replace logarithms by $\log(s+16)$ and
$\log\log(s+16)$; this does not change the asymptotics.

\begin{corollary}\label{cor:time}
Using the cited manuscript theorem, $|\MT(H)|$ and its entire cardinality
distribution are computable deterministically in
\begin{equation}\label{eq:time}
 (s+I+1)^{O(1)}
 2^{O(s\log\log(s+16)/\sqrt{\log(s+16)})}
\end{equation}
time and polynomial space. If only the documented Lean scope is used, replace
the numerator's $\log\log$ by $(\log\log)^2$.
If $I+J=O(s)$, the manuscript yields
\[
 (s+I+1)^{O(1)}2^{O(s/\sqrt{\log(s+16)})}.
\]
In particular, the last bound applies to hypergraphs of fixed maximum edge
size, with constants depending on that size.
\end{corollary}
\begin{proof}
We have $N=O(s)$ and $M=O(s+I+J)$, while $J\le pI/2$. Ordinary binary-index
encodings have polynomial bit length. The canonical encoding described in
the Lean comparator uses unary indices and still has polynomial bit length
for these formulas. Thus all input factors remain polynomial. For a fixed
cardinality the additional variables and clauses are $O(p)$; polynomially
many sequential calls preserve the bound and polynomial space. If
$I+J=O(s)$ then $M=O(N)$, including the size circuit. Each oracle answer has
at most $N+1$ bits, so exact subtraction requires only polynomial resources.
\end{proof}

\subsection{Minimal dominating sets and upper domination}
For a simple undirected graph $G=(V,E)$, use the indexed hypergraph
$(V,(N[v])_{v\in V})$. A set is a minimal transversal exactly when it is a
minimal dominating set~\cite{bbcw}. For $n=|V|\ge1$, we obtain
\[
 N=5n-1,\qquad
 M=10n+2|E|-3+\sum_{v\in V}\binom{\deg(v)+1}{2}.
\]
The expression for $M$ follows by substituting $p=q=n$ and
$I=n+2|E|$ in Theorem~\ref{thm:main}. Even for dense graphs it is $O(n^3)$,
and the running time in \eqref{eq:time} depends exponentially only on $n$.
Fixed maximum degree gives the improved exponent $O(n/\sqrt{\log n})$.

Upper domination asks for the largest cardinality of a minimal dominating
set~\cite{upper}. Computing the distribution and taking its largest nonzero
index solves it within the same time bound. A graph always has a minimal
dominating set, including the empty graph under the empty-set convention.

\subsection{Minimal set covers and other direct applications}
Given $a$ indexed selectable sets over $b$ universe elements, take one
hypergraph vertex per selectable set and one edge per element, consisting
of the sets that contain it. Minimal covers correspond bijectively to
minimal transversals. The parameter is $a+b$, not the number of incidences;
duplicates among selectable sets remain distinct choices. An uncovered
element gives an empty edge and zero covers.

The same statement counts minimal satisfying assignments of a monotone CNF,
using its variables as vertices and clauses as edges. Open neighborhoods
instead of closed neighborhoods give minimal total dominating sets; a graph
with an isolated vertex then has count zero. Complementation gives maximal
independent vertex sets of a hypergraph, where independence means containing
no entire edge. These are consequences of the same reduction, not separate
core constructions.

\subsection{Witnesses and exact sampling}
Because Lemma~\ref{lem:pointwise} holds for each $D$, adding units on $x$
gives exact nonnegative conditional counts even though the internal counts
are subtracted. If a target count is positive, fix the selection variables
successively to a branch with positive conditional count. This recovers one
minimal transversal of a specified feasible size in $O(p)$ additional pairs
of calls. It also recovers an upper dominating set after finding its size.
Selecting each branch with probability equal to its conditional count divided
by the parent count gives an exactly uniform sample. Exact integer rejection
sampling supplies these probabilities with polynomial expected overhead.
These are algorithms for objects represented by the $x$ bits, not uniform
sampling of the auxiliary SAT assignments.

\section{Prior work and claim discipline}\label{sec:prior}
The private-edge characterization, the neighborhood-hypergraph correspondence,
inclusion--exclusion, and uniquely extendible gate encodings are established
ingredients. Prior negative-weight model-counting transformations specifically
eliminate existential witnesses by cancellation~\cite{skolem}. It would be
incorrect to claim that the general cancellation idea was invented here.

The supplied related preprint~\cite{claude} already handles ordinary dominating
set counting, including cardinality restrictions, in Corollary~7. Minimality
is an additional condition and cannot be obtained merely by renaming that
result. The present note supplies an explicit construction and its variable
accounting. Moreover, the existing negative-weight technique already gives
the same asymptotic linear-variable existence result by the specialization
below. Thus that existence result alone is not a defensible methodological
novelty claim. Whether this exact application or its tighter constant was
published previously remains unconfirmed; search failure is not a priority proof.

\subsection{An explicit specialization of the earlier technique}
Here is our parameter accounting for a baseline obtained by specializing the
negative-weight rules of Definition~3 and Proposition~5 in~\cite{skolem}.
This derivation makes the relationship checkable rather than merely comparing
keywords. Use variables $x_v,u_v$ and $b_i,c_i,t_i$, giving $2p+3q$ variables.
Assign weight $-1$ to a false $u_v,c_i,t_i$ and weight $+1$ to every other
variable-value pair. Include
\begin{align*}
 &(b_i\vee\neg x_v\vee\neg x_w),\quad
   (c_i\vee\neg x_v\vee\neg x_w)
   &&(\{v,w\}\subseteq e_i),\\
 &(c_i\vee b_i) &&(i\in[q]),\\
 &(t_i\vee\neg x_v) &&(v\in e_i),\\
 &(u_v\vee x_v) &&(v\in U),\\
 &(u_v\vee b_i) &&(v\in e_i).
\end{align*}
Fix $x$. Summing over $c_i$ enforces, with weight one,
$b_i=\ind{|e_i\cap D|\ge2}$: if a pair is selected both $b_i,c_i$ are forced
true; if no pair is selected, false $b_i$ forces true $c_i$, whereas true
$b_i$ leaves $c_i$ free and cancels. Summing over $t_i$ enforces that $e_i$
is hit. Once all edges are hit and $b$ has the indicated values, summing over
$u_v$ enforces $\neg x_v\vee\bigvee_{i:v\in e_i}\neg b_i$, which is exactly
private-edge availability for each selected vertex. These filters contribute
one precisely for minimal transversals and zero otherwise.

There are $p+2q$ signed bits. Their number of false values has the parity of
their number of true values plus the fixed constant $p+2q$. A parity chain
therefore converts this baseline to two ordinary counts with
$3p+5q-1$ variables for $p+q\ge1$, merely swapping the two counts when that
constant is odd. Theorem~\ref{thm:main} saves $3q$ variables relative to this
particular baseline. Neither construction is asserted optimal. This comparison
supports an expository or application contribution, not a claim that prior
methods cannot obtain a linear-variable reduction.

For beta-acyclic hypergraphs, polynomial-time minimal transversal counting
was already known; the neighborhood connection also gives results on strongly
chordal graphs~\cite{bbcw}. Those tractable subclasses do not imply the present
arbitrary-incidence bound. Conversely, a general subexponential bound is not
a claim to improve those polynomial-time algorithms.

The deliverable is a proved and reproducible reduction with explicit
algorithmic corollaries. Its status is a research note, not a claim of
publication-ready novelty or a new proof of the external ETH result.

\section{Verification and reproducibility}\label{sec:verify}
The accompanying standard-library Python implementation is in
\texttt{code/}. It generates the actual two CNFs, optionally adds the balanced
cardinality circuit, and exports ordinary DIMACS. Its exact DPLL counter is
a small-instance verification oracle, not an implementation of the cited
subexponential algorithm. Unit propagation and free-variable factors must
preserve declared-variable semantics.

Validation compares three independently expressed objects: direct enumeration
of minimal transversals; the signed identity, with the $P$ sum factored; and,
on smaller cases, the difference of actual generated CNF model counts.
The verification suite checks all 278 simple hypergraphs on zero through three
vertices using actual CNF counts, both in total and by size (2,720 oracle
calls in that exhaustive section). All 65,536 simple hypergraphs on four
vertices are checked against both factored signed identities, not exhaustively
through the CNF oracle. It also covers gate truth tables, unique extension,
random CNFs for the DPLL oracle, repeated-edge cases, all 76 simple graphs on
zero through four vertices, and set-cover examples. The prior-technique
baseline is independently checked on all 22 simple hypergraphs with at most
two vertices, including all 74 fixed-$D$ contributions (192 oracle calls).
Running \texttt{code/verify.py} generates a report with coverage and running times
under \texttt{workspace/}.
The proof above, not a finite test suite, establishes the general result.
The local reduction itself has not been formalized in Lean.

\section{Limits of the two optional directions}\label{sec:limits}
For connected dominating sets, a consistent-cut filter with a fixed root
counts $2^{c-1}$ cuts when there are $c$ components~\cite{cutcount}. This
removes disconnected objects modulo two but leaves positive multiplicities
over the integers. A spanning-tree certificate similarly weights a set by
its number of spanning trees. Neither observation supplies a one-per-object
integer count. The present private-edge argument relies on monotonicity and
independent local witnesses; it supplies no such connectivity filter.
The separate note \texttt{papers/connected\_domination.tex} instead proves a
$3n$-variable parity representation, a polynomial-marker extension of the
external interpolation algorithm, and randomized subexponential decision,
search, and minimum-size optimization. Its parity argument does not resolve
the exact integer counting question discussed in this paragraph.

For bounded-degree feedback vertex set, a graph is acyclic precisely when
its retained vertices meet no cycle in full. Turning every cycle into a
hyperedge can produce exponentially many edges even at bounded degree.
The parameter $p+q$ would then be too large, and our theorem gives no desired
bound. An ordering encoding with $O(n\log n)$ bits also does not suffice:
substitution into the imported exponent yields order
$n\sqrt{\log n}\log\log n$, which is not $o(n)$.
No low-variable fixed-width reduction for this direction is established here.
These are remaining mathematical problems, not implementation tasks already
completed by the main theorem.

Finally, fast counting does not imply subexponential total-time enumeration.
A disjoint union of triangles has $3^{n/3}$ minimal dominating sets (one
vertex per triangle), so simply writing them all already takes exponential
time. The sampling and single-witness procedures do not contradict this.

\begin{thebibliography}{9}
\bibitem{oai} OpenAI.
\emph{The Exponential Time Hypothesis is False}.
17 October 2026, revised 25 October 2026, main theorem.
\href{https://github.com/openai/math/blob/main/preprints/The-Exponential-Time-Hypothesis-is-False-October-25-2026/paper.pdf}{Manuscript};
\href{https://raw.githubusercontent.com/openai/math/refs/heads/main/preprints/The-Exponential-Time-Hypothesis-is-False-October-25-2026/build/main_full_version.tex}{source}.

\bibitem{scope} OpenAI.
\emph{Lean formalization scope, problem 378}.
Accessed 29 October 2026.
\href{https://github.com/openai/math/blob/main/lean/docs/378.md}{Scope document};
\href{https://github.com/openai/math/blob/main/lean/ComparatorChallenges/SubexponentialSAT.lean}{comparator statement}.

\bibitem{claude} Claude / Anthropic.
\emph{A $2^{o(n)}$ Time Algorithm for 3-SAT}.
26 October 2026, preprint 2610.00577v1, Corollary 7, pp.21--22.
\href{https://hexagonmath.org/pdf/2610.00577v1}{Primary manuscript}.

\bibitem{bbcw} B. Bergougnoux, F. Capelli, and M. Kant\'e.
\emph{Counting Minimal Transversals of Beta-Acyclic Hypergraphs}.
\href{https://arxiv.org/abs/1808.05017}{arXiv:1808.05017}.

\bibitem{skolem} G. Van den Broeck, W. Meert, and A. Darwiche.
\emph{Skolemization for Weighted First-Order Model Counting}.
KR 2014.
\href{https://arxiv.org/abs/1312.5378}{arXiv:1312.5378}.

\bibitem{upper} C. Bazgan et al.
\emph{Algorithmic Aspects of Upper Domination}.
\href{https://arxiv.org/abs/1506.07260}{arXiv:1506.07260}.

\bibitem{cutcount} M. Cygan et al.
\emph{Solving connectivity problems parameterized by treewidth in single
exponential time}.
\href{https://arxiv.org/abs/1103.0534}{arXiv:1103.0534}.
\end{thebibliography}
\end{document}
